% ===========================================================================
% Web preamble (article target for LaTeXML + BookML) — loaded by drivers/web.tex.
%
% CLASSIC TeX BY CONSTRUCTION: no xparse, no expl3, no engine-specific escape
% hatches.  This is the collapse of the donor course's beamer->article shim
% layer (PHYS 567 preamble_html.tex + phys-lxml-compat.tex) into one ordinary
% preamble — LaTeXML 0.8.8 parses it directly, and it also compiles under
% plain pdflatex (which BookML's build relies on: the PDF pass runs first and
% doubles as a regression check).
%
% Builds run from the REPO ROOT: all paths are root-relative.
% ===========================================================================
\usepackage[utf8]{inputenc}
\usepackage{amsmath,amssymb,mathtools}
\usepackage{amsthm}
\usepackage{array}
\usepackage{booktabs}
\usepackage{graphicx}
\usepackage{xcolor}
\usepackage{hyperref}
\usepackage{tikz}
\usetikzlibrary{arrows.meta,calc,positioning,shapes.geometric}
% overlay-beamer-styles is beamer-only; make its keys harmless no-ops here.
\tikzset{visible on/.style={}}

% --- Kit palette (keep in sync with preamble-slides.tex and style.css) -------
\definecolor{physblue}{HTML}{13294B}
\definecolor{physorange}{HTML}{C84113}
\definecolor{physgray}{HTML}{F3F4F6}
\definecolor{physink}{HTML}{111827}
\definecolor{physgreen}{HTML}{1B7A3D}

% --- BookML (when present) ---------------------------------------------------
% nomathjax: native MathML, zero JS (the site's stance).  style=plain: no
% gitbook chrome.  tikzpicture environments render as dvisvgm SVG images —
% the accessible figure path (donor pilot, streamlining.md §8.6).
% Guarded so the document still compiles under plain pdflatex without the
% bookml/ directory (debug builds outside the BookML flow).
\IfFileExists{bookml/bookml.sty}{%
  \usepackage[nomathjax,style=plain]{bookml/bookml}%
  \bmlImageEnvironment{tikzpicture}%
}{}

% --- Lecture title metadata --------------------------------------------------
\newcommand{\kitlecturetitle}[1]{%
  \title{#1}\author{\coursecode}\date{\courseterm}}

% ===========================================================================
% Beamer -> article shims (all classic TeX)
% ===========================================================================
\makeatletter

% --- helpers: strip one optional <...>; keep/gobble a following {group} ------
\def\kit@ang#1{\@ifnextchar<{\kit@ang@{#1}}{#1}}
\def\kit@ang@#1<#2>{#1}
\long\def\kitunwrapgroup#1{#1}
\long\def\kitgobblegroup#1{}
\def\kit@maybegroup{\@ifnextchar\bgroup{\kitunwrapgroup}{}}

% --- title page: beamer's \titlepage inside the title frame -> \maketitle ----
% AtBeginDocument, not a preamble-time \def: LaTeXML re-establishes its own
% \titlepage constructor after the preamble runs, so a plain override here is
% silently lost (donor pilot lesson, 2026-08-02).  Harmless under pdflatex.
\AtBeginDocument{\def\titlepage{\maketitle}\def\endtitlepage{}}

% --- frame -> subsection heading (final, fully-revealed state) ---------------
% \section stays article \section (=> h2); each frame title becomes a
% \subsection (=> h3, and a real <section> element under LaTeXML).  A
% title-less frame (the title page) emits no heading.
\def\frame{\frame@scan}
\newcommand\frame@scan[1][]{\@ifnextchar\bgroup{\frame@titled}{}}
\newcommand\frame@titled[1]{\subsection{#1}}
\def\endframe{}
\def\frametitle{\kit@ang\kit@frametitle}
\def\kit@frametitle#1{\subsection{#1}}
\def\framesubtitle{\kit@ang\kit@framesubtitle}
\def\kit@framesubtitle#1{\par\noindent\textbf{#1}\par}
% Frames are unnumbered in the archival document; \sections keep 1..N.
\setcounter{secnumdepth}{1}

% --- overlays: drop the <...> spec, keep the final-state content -------------
\def\onslide{\kit@ang\kit@maybegroup}
\let\uncover\onslide
\let\visible\onslide
\def\pause{\kit@ang\relax}
% \alt<ov>{shown}{hidden} -> shown (also what TikZ's \draw<1-> etc. expand to)
\def\alt{\kit@ang\kit@alttake}
\long\def\kit@alttake#1#2{#1}
% \only: a CLOSED range (<1>, <1-2>) is staging the final state does NOT show —
% drop the group; open ranges (<2->) and rangeless \only keep it.
% Ends-with-dash test via the delimited-match idiom.
\def\only{\@ifnextchar<{\kit@onlyspec}{\kit@maybegroup}}
\def\kit@onlyspec<#1>{\kit@endsdash#1\kit@mk-\kit@mk\kit@stop}
\long\def\kit@endsdash#1-\kit@mk#2\kit@stop{%
  \if\relax\detokenize{#2}\relax
    \expandafter\kit@dropgroup
  \else
    \expandafter\kit@maybegroup
  \fi}
\def\kit@dropgroup{\@ifnextchar\bgroup{\kitgobblegroup}{}}

% --- \item<ov> ---------------------------------------------------------------
\let\kit@origitem\item
\def\item{\@ifnextchar<{\kit@itemang}{\kit@origitem}}
\def\kit@itemang<#1>{\kit@origitem}

% --- lists may carry <overlay> and/or [template] after \begin{...} -----------
\let\kit@origitemize\itemize         \let\kit@endorigitemize\enditemize
\let\kit@origenumerate\enumerate     \let\kit@endorigenumerate\endenumerate
\let\kit@origdescription\description \let\kit@endorigdescription\enddescription
\def\kit@liststrip#1{\@ifnextchar<{\kit@liststrip@ang{#1}}{\kit@liststrip@opt{#1}}}
\def\kit@liststrip@ang#1<#2>{\kit@liststrip@opt{#1}}
\def\kit@liststrip@opt#1{\@ifnextchar[{\kit@liststrip@opt@{#1}}{#1}}
\def\kit@liststrip@opt@#1[#2]{#1}
\renewenvironment{itemize}{\kit@liststrip\kit@origitemize}{\kit@endorigitemize}
\renewenvironment{enumerate}{\kit@liststrip\kit@origenumerate}{\kit@endorigenumerate}
\renewenvironment{description}{\kit@liststrip\kit@origdescription}{\kit@endorigdescription}

% --- proof: strip a leading <overlay>, keep amsthm's optional [name] ---------
\let\kit@origproof\proof  \let\kit@endorigproof\endproof
\def\kit@proofstrip{\@ifnextchar<{\kit@proofang}{\kit@proofopt}}
\def\kit@proofang<#1>{\kit@proofopt}
\def\kit@proofopt{\@ifnextchar[{\kit@proofname}{\kit@origproof}}
\def\kit@proofname[#1]{\kit@origproof[#1]}
\renewenvironment{proof}{\kit@proofstrip}{\kit@endorigproof}

% --- \alert -> bold + physorange ---------------------------------------------
\def\alert{\kit@ang\kit@alertdo}
\long\def\kit@alertdo#1{\textbf{\textcolor{physorange}{#1}}}

% --- columns / column -> linear top-to-bottom stacking (= reading order) -----
\newenvironment{columns}{\kit@colsbegin}{\par}
\def\kit@colsbegin{\@ifnextchar[{\kit@colsopt}{\par}}
\def\kit@colsopt[#1]{\par}
\newenvironment{column}{\kit@colbegin}{\par}
\def\kit@colbegin{\@ifnextchar[{\kit@colopt}{\kit@colwidth}}
\def\kit@colopt[#1]{\kit@colwidth}
\def\kit@colwidth#1{\par}

% --- block / alertblock / exampleblock ---------------------------------------
% Rendered as bold-title paragraphs for now; a LaTeXML binding giving semantic
% <div class="block"> containers (for card styling) is a planned enhancement.
% Overlay specs are stripped on BOTH sides of the title.
\def\kit@blockbegin{\@ifnextchar<{\kit@blockpre}{\kit@blocktit}}
\def\kit@blockpre<#1>{\kit@blocktit}
\def\kit@blocktit#1{\par\medskip\noindent{\bfseries #1}\par\nobreak
  \@ifnextchar<{\kit@blockpost}{}}
\def\kit@blockpost<#1>{}
\newenvironment{block}{\kit@blockbegin}{\par\medskip}
\newenvironment{alertblock}{\kit@blockbegin}{\par\medskip}
\newenvironment{exampleblock}{\kit@blockbegin}{\par\medskip}

% --- \case{N} and \textcircled safety ----------------------------------------
% LaTeXML stringifies \textcircled's argument (decorated arguments come out as
% literal TeX — donor lesson 2026-08-02), and a circled digit conveys nothing a
% screen reader can't get from the bare number: parenthesized form on the web.
\newcommand{\case}[1]{(#1)}
\def\textcircled#1{(#1)}

% --- figure alt-text registry ------------------------------------------------
% Authored alt text lives in alt/lectureNN.tex sidecars (keyed
% "lectureNN/figure-name" by OWNING lecture; never regenerated).  \figalt must
% stay expandable (plain \newcommand) so it can be materialized into the
% output attribute.
\newcommand{\setfigalt}[2]{\expandafter\def\csname kitalt@#1\endcsname{#2}}
\newcommand{\figalt}[1]{%
  \ifcsname kitalt@#1\endcsname\csname kitalt@#1\endcsname\else[ALT MISSING: #1]\fi}

% --- shared-figure alt borrowing ---------------------------------------------
% A lecture reusing another lecture's figure declares \usealtfrom{lectureMM};
% the owner's sidecar is loaded so its keys resolve here.
\newcommand{\usealtfrom}[1]{\IfFileExists{alt/#1.tex}{\input{alt/#1.tex}}{}}

% --- \fig{path}{alt-key} -----------------------------------------------------
% The TikZ source is \input (BookML images it via dvisvgm); \kitfigdesc is the
% hook that attaches \figalt{#2} to the generated <img alt> — a no-op until the
% LaTeXML binding lands (kit work item A3), harmless under pdflatex.
\newcommand{\kitfigdesc}[1]{}
\newcommand{\fig}[2]{%
  \par\begin{center}\kitfigdesc{\figalt{#2}}\input{#1}\end{center}\par}

% --- announcements: dropped by construction ----------------------------------
% The archive is announcement-free; course.py's --keep-announcements defines
% \kitkeepannouncements to flip them on for a one-off build.
\ifdefined\kitkeepannouncements
  \newcommand{\kitannouncementspath}{announcements/lecture\kitlecture.tex}
  \newcommand{\announcementsframe}{%
    \IfFileExists{\kitannouncementspath}%
      {\begin{frame}{Announcements}\input{\kitannouncementspath}\end{frame}}{}}
  \newcommand{\announcementsblock}{%
    \IfFileExists{\kitannouncementspath}%
      {\begin{block}{Announcements}\input{\kitannouncementspath}\end{block}}{}}
\else
  \newcommand{\announcementsframe}{}
  \newcommand{\announcementsblock}{}
\fi

\makeatother
